%
%  FMICS 2026 Tool Demonstration Paper
%  stC: A C-Syntax Language for Formal Verification of Safety-Critical Embedded Systems
%
\documentclass[runningheads]{llncs}

\usepackage[T1]{fontenc}
\usepackage{graphicx}
\usepackage{listings}
\usepackage{xcolor}
\usepackage{hyperref}
\usepackage{amsmath}

% stC code listing style
\lstdefinelanguage{stC}{
  keywords={module, import, type, subtype, const, foreign, private, volatile,
            struct, union, enum, range, modular, float, fixed, decimal,
            digits, delta, if, else, while, for, do, switch, case, default,
            break, continue, return, goto, in, out, inout, void, bool, char,
            true, false, pre, post, global, reads, writes,
            loop_invariant, loop_variant, type_invariant, assert,
            offset, bit, bits, at, sizeof},
  sensitive=true,
  morecomment=[l]{//},
  morecomment=[s]{/*}{*/},
  morestring=[b]",
  morestring=[b]'
}

\lstset{
  language=stC,
  basicstyle=\ttfamily\small,
  keywordstyle=\bfseries,
  commentstyle=\itshape\color{gray},
  stringstyle=\color{black},
  numbers=left,
  numberstyle=\tiny,
  numbersep=5pt,
  frame=single,
  breaklines=true,
  captionpos=b,
  tabsize=2
}

\lstdefinelanguage{AdaCode}{
  keywords={type, is, range, mod, record, end, with, Address, Volatile,
            System, Storage_Elements, To_Address, at, for, use},
  sensitive=true,
  morecomment=[l]{--},
  morestring=[b]"
}

\begin{document}

\title{stC: A C-Syntax Language for Formal Verification\\of Safety-Critical Embedded Systems}

\author{Germ\'an Rivera\inst{1}}

\authorrunning{G. Rivera}

\institute{%
  \email{jgrivera67@gmail.com}\\
  \url{https://github.com/jgrivera67/stcc}
}

\maketitle

\begin{abstract}
Safety-critical embedded software requires both strong correctness guarantees and
efficient low-level hardware control. Ada/SPARK provides excellent support for
formal verification but its syntax discourages adoption among C-fluent embedded
engineers. We present \emph{stC} (\emph{Strongly-Typed C}), a new systems programming
language that combines C-like syntax with Ada-level strong typing and formal
verification support. stC is transpiled to SPARK~Ada, enabling the full
SPARK verification toolchain on code written in a familiar C-like notation.
Key features include: range and modular integer types that eliminate overflow
errors, pre/post-condition contracts, global effect clauses, explicit
memory-mapped I/O register layout with bit-precise field annotations,
and pointerless hardware access. We describe the \texttt{stcc} transpiler---itself
written in SPARK Ada---and its VSCode language-server extension, and
illustrate the workflow on a GPIO driver case study.
\end{abstract}

\keywords{systems programming \and formal verification \and embedded systems
  \and transpiler \and Ada/SPARK \and strong typing \and MMIO \and contracts}

%------------------------------------------------------------------
\section{Introduction}
\label{sec:intro}
%------------------------------------------------------------------

Most safety-critical embedded software is written in C. Despite C's dominance, the
language is notoriously difficult to verify formally: implicit integer promotions,
unconstrained pointer arithmetic, undefined behaviour, and the absence of built-in
contracts place a heavy burden on verification tools that must reconstruct
programmer intent from an underspecified semantics~\cite{misra-c}.

Ada/SPARK~\cite{SparkAda} addresses these problems head-on. Its strong type system
eliminates implicit conversions, its contract language expresses pre- and
post-conditions directly in source code, and the SPARK subset admits formal proof
via SMT solvers. Yet industrial adoption of Ada/SPARK remains limited among
engineers with a C background, largely because of Ada's verbose syntax and
unfamiliar programming idioms~\cite{ada-industry}.

We propose \emph{stC} (\emph{Strongly-Typed C}) as a bridge. stC is a new
programming language for safety-critical embedded systems that:
\begin{itemize}
  \item uses familiar C-like syntax so C engineers can be productive immediately;
  \item enforces Ada-level strong typing---no implicit conversions, all integer
        types declared with explicit bounds;
  \item provides Ada/SPARK-style contracts (pre/post conditions, global effect
        clauses, loop invariants) directly in source;
  \item offers precise bit-level control over hardware register layouts for
        memory-mapped I/O (MMIO), including pointerless variable placement
        at fixed addresses;
  \item transpiles to \emph{SPARK Ada 2022}, making the full SPARK proof
        toolchain available to stC programs.
\end{itemize}

The \texttt{stcc} compiler is itself written in SPARK Ada, serving as a
demonstration that the formal methods it enables are practical for real
tool development. The tool is open-source and available at
\url{https://github.com/jgrivera67/stcc}.

%------------------------------------------------------------------
\section{The stC Language}
\label{sec:language}
%------------------------------------------------------------------

\subsection{Design Principles}

stC follows four design principles:

\noindent\textbf{C syntax, Ada semantics.}
The language looks like C~--- curly braces, semicolons, \texttt{if}/\texttt{while}/\texttt{for} ---
but its semantics follow Ada. There are no header files, no preprocessor,
no implicit type conversions, and no built-in integer types.

\noindent\textbf{Declarations before use is optional.}
A module's declarations may appear in any order; there is no need for forward
declarations or function prototypes. The compiler resolves all references in a
single-pass over the module graph.

\noindent\textbf{Pointers are the exception.}
Array parameters carry their length. Record field access and MMIO register
access require no pointer arithmetic. Pointers exist but are not the default
mechanism for indirection.

\noindent\textbf{Verification readiness by construction.}
Type bounds, contracts, and effect annotations are first-class syntax, not
comments or annotations that can be stripped. Generated SPARK Ada preserves
all these guarantees.

\subsection{Modules}

A stC source file declares a named module and may import other modules:

\begin{lstlisting}[caption={Module structure}]
module gpio.driver {
  import hardware.types;
  import platform.registers as regs;

  // Public declarations ...

  private
  // Private declarations ...
}
\end{lstlisting}

Hierarchical module names (\texttt{parent.child}) correspond to Ada child
packages, giving the module system a familiar structure for SPARK users.

%------------------------------------------------------------------
\section{Type System}
\label{sec:types}
%------------------------------------------------------------------

stC has no built-in integer types. Every integer type must be explicitly declared
with either a signed range or a modular (unsigned wraparound) specification:

\begin{lstlisting}[caption={Scalar type declarations}]
// Signed range types -- map to Ada range types
type range 0 .. 255    byte_t;
type range -32768 .. 32767 int16_t;
type range 0 .. 7      uint3_t;   // 3-bit logic

// Modular (unsigned) types -- wrap on overflow
type modular 256         uint8_t;
type modular 1 << 16     uint16_t;
type modular 1 << 32     uint32_t;
type modular 1 << machine_width  address_t;
\end{lstlisting}

Range types map to Ada \texttt{range} types, while modular types map to Ada
\texttt{mod} types. Both carry value constraints that are enforced statically
by SPARK and at runtime by Ada range checks.

\subsection{Floating-Point and Fixed-Point Types}

stC follows Ada's model for real types, distinguishing floating-point precision
from fixed-point resolution:

\begin{lstlisting}[caption={Real number type declarations}]
// IEEE floating-point with precision guarantee
type float digits 6   float32_t;
type float digits 15  float64_t;

// Sensor type: 0.001 resolution, physical range
type float digits 6 range -273.15 .. 1000.0  temperature_t;

// Fixed-point: exact delta, useful for control loops
type fixed delta 0.001 range 0.0 .. 5.0  voltage_t;  // mV resolution

// Decimal: base-10 exact, useful for financial values
type decimal delta 0.01 digits 10  money_t;
\end{lstlisting}

\subsection{Subtypes}

A subtype constrains an existing type to a narrower range:

\begin{lstlisting}[caption={Subtype declaration}]
type range 0 .. 255 byte_t;
subtype byte_t range 0 .. 127 ascii_t;
\end{lstlisting}

Generated Ada preserves the subtype relationship, allowing SPARK to verify
that values remain within bounds through all assignments.

\subsection{Composite Types}

Structs, unions, and enumerations use C-like syntax:

\begin{lstlisting}[caption={Composite types}]
type enum { idle, running, stopped, error }  status_t;

type struct {
  byte_t x;
  byte_t y;
} point_t;

type union {
  uint32_t word;
  uint8_t  bytes[4];
} raw_data_t;
\end{lstlisting}

%------------------------------------------------------------------
\section{Contracts and Verification Support}
\label{sec:contracts}
%------------------------------------------------------------------

stC supports Ada/SPARK-style contracts as first-class syntax using
C23-style attribute syntax (\texttt{[[...]]}).

\begin{lstlisting}[caption={Pre/post-condition contracts}]
byte_t saturating_add(in byte_t a, in byte_t b)
  [[pre(a <= 255 - b)]]
  [[post(result == a + b)]]
{
  return a + b;
}
\end{lstlisting}

\noindent\textbf{Global effect clauses} declare which global variables a
function reads or writes. This information is indispensable for SPARK's
information-flow analysis:

\begin{lstlisting}[caption={Global effect and contracts}]
uint32_t increment_counter()
  [[global(reads => (global_counter), writes => (global_sum))]]
  [[pre(global_counter < 1000)]]
  [[post(global_sum > 0)]]
{
  global_sum = global_sum + global_counter;
  return global_sum;
}
\end{lstlisting}

\noindent\textbf{Loop contracts} support termination and invariant proofs:

\begin{lstlisting}[caption={Loop invariant and variant}]
while (i <= index_t'Last)
  [[loop_invariant(sum >= 0 && i >= index_t'First)]]
  [[loop_variant(index_t'Last - i)]]
{
  sum = sum + arr[i];
  i = i + 1;
}
\end{lstlisting}

\noindent\textbf{Parameter modes} make data flow explicit, aiding both human
readers and verification tools:

\begin{lstlisting}[caption={Parameter modes}]
void compute(in sensor_t reading,
             out result_t result,
             inout state_t state);
\end{lstlisting}

\noindent\textbf{Type attributes} follow Ada's \texttt{'First}, \texttt{'Last},
\texttt{'Range}, and \texttt{'Size} notation, allowing range-safe loop bounds
and size assertions to be written without magic numbers:

\begin{lstlisting}
i = index_t'First;
while (i <= index_t'Last) { ... }
\end{lstlisting}

All contracts are forwarded verbatim into the generated SPARK Ada, where the
GNATprove tool can discharge them via SMT solvers or flag unverified
obligations.

%------------------------------------------------------------------
\section{Memory-Mapped I/O and Embedded Hardware Support}
\label{sec:mmio}
%------------------------------------------------------------------

MMIO register programming is a domain where C falls short: bit-field layout
is compiler-defined, volatile semantics are underspecified, and address binding
requires unsafe pointer casts. stC provides explicit, formally grounded alternatives.

\subsection{Bit-Field Annotations}

Struct fields may carry \texttt{[[offset(byte) bit(n)]]} or
\texttt{[[offset(byte) bits(start..end)]]} annotations that precisely specify
each field's byte offset and bit position within the register:

\begin{lstlisting}[caption={GPIO 32-bit control register with bit fields}]
type struct {
  bool    enable        [[offset(0) bit(0)]];   // Byte 0, bit 0
  bool    irq_enable    [[offset(0) bit(1)]];   // Byte 0, bit 1
  uint3_t mode          [[offset(0) bits(2..4)]]; // Byte 0, bits 2-4
  uint8_t reserved      [[offset(1) bits(0..7)]]; // Byte 1 (reserved)
  uint8_t pin_value     [[offset(2) bits(0..7)]]; // Byte 2: pin data
  uint8_t irq_status    [[offset(3) bits(0..7)]]; // Byte 3: interrupts
} gpio_control_t;
\end{lstlisting}

An \emph{all-or-nothing} rule applies: either every field in a struct carries
offset/bit annotations, or none do. This prevents accidental mixing of
compiler-managed and programmer-managed layout---a common source of subtle
bugs in embedded C.

These annotations translate to Ada representation clauses:

\begin{lstlisting}[language=AdaCode, caption={Generated SPARK Ada representation clause}]
for GPIO_Control_T use record
   Enable     at 0 range 0 .. 0;
   Irq_Enable at 0 range 1 .. 1;
   Mode       at 0 range 2 .. 4;
   Reserved   at 1 range 0 .. 7;
   Pin_Value  at 2 range 0 .. 7;
   Irq_Status at 3 range 0 .. 7;
end record;
\end{lstlisting}

\subsection{Pointerless MMIO Variable Placement}

Ada's \texttt{Address} aspect allows a variable to be placed at a specific
memory address without pointer arithmetic. stC exposes this via the
\texttt{[[at(address)]]} variable attribute, combined with the \texttt{volatile}
qualifier:

\begin{lstlisting}[caption={Pointerless MMIO register access}]
// Hardware registers placed at fixed addresses -- no pointers needed
volatile gpio_control_t GPIO_PORT_A [[at(0x40000000)]];
volatile gpio_control_t GPIO_PORT_B [[at(0x40000100)]];

void init_gpio() {
  // Direct field access -- no dereference operator
  GPIO_PORT_A.enable     = true;
  GPIO_PORT_A.irq_enable = true;
  GPIO_PORT_A.mode       = 3;
}
\end{lstlisting}

This maps to SPARK Ada with the \texttt{Address} and \texttt{Volatile} aspects:

\begin{lstlisting}[language=AdaCode, caption={Generated SPARK Ada for MMIO variable}]
GPIO_Port_A : GPIO_Control_T
   with Address  =>
        System.Storage_Elements.To_Address (16#40000000#),
        Volatile;
\end{lstlisting}

Pointerless MMIO access eliminates an entire class of unsafe operations
(pointer casts, volatile pointer misuse) while generating code that is
formally verifiable with SPARK's information-flow analysis.

%------------------------------------------------------------------
\section{The \texttt{stcc} Transpiler}
\label{sec:transpiler}
%------------------------------------------------------------------

\texttt{stcc} is a source-to-source compiler (transpiler) that translates stC
to SPARK Ada 2022. It is written entirely in SPARK Ada and builds
with zero warnings under the GNAT compiler with full restriction pragmas.

\subsection{Architecture}

The compiler pipeline consists of three stages:
\begin{enumerate}
  \item \textbf{Lexer}: Tokenises the stC source. Implements all stC keywords,
        numeric literals in decimal, hexadecimal (0x), and binary (0b) notation,
        digit separators (\texttt{\_} and \texttt{'}), and floating-point
        literals.
  \item \textbf{Parser}: Builds an abstract syntax tree (AST). A recursive-descent
        parser handles declarations, statements, expressions, type annotations,
        contracts, and attribute blocks.
  \item \textbf{Code Generator}: Traverses the AST and emits SPARK Ada 2022
        source. Original stC line numbers are preserved as comments in the
        generated code to aid debugging.
\end{enumerate}

\subsection{Build System}

The project uses the Alire~\cite{alire} package manager, making the tool
easy to install on Linux and macOS:

\begin{lstlisting}[language=sh, numbers=none]
$ alr get stcc && alr build
$ stcc my_module.stc
\end{lstlisting}

%------------------------------------------------------------------
\section{Tool Demonstration}
\label{sec:demo}
%------------------------------------------------------------------

\subsection{VSCode Language Extension}

stC is accompanied by a VSCode extension providing a full Language Server
Protocol (LSP) implementation~\cite{lsp}. The extension offers:

\begin{itemize}
  \item \textbf{Syntax highlighting} for all stC constructs including
        contract keywords (\texttt{pre}, \texttt{post}, \texttt{global}),
        type constructors (\texttt{range}, \texttt{modular}), hardware
        attribute keywords (\texttt{offset}, \texttt{bit}, \texttt{bits},
        \texttt{at}), and parameter modes (\texttt{in}, \texttt{out},
        \texttt{inout}).
  \item \textbf{Real-time diagnostics}: the extension invokes \texttt{stcc}
        on the document as the user types (with configurable debounce delay)
        and displays compiler errors inline with exact line and column
        positions.
  \item \textbf{Go-to-definition}: navigates to type and function definitions
        within the module.
  \item \textbf{Document symbols}: the symbol outline lists all types,
        functions, and subtypes for rapid navigation.
  \item \textbf{Code completion}: suggests all stC keywords and known
        user-defined type identifiers.
\end{itemize}

\subsection{Demo Scenario: GPIO Driver}

We demonstrate the complete workflow for a GPIO peripheral driver. Starting
from a blank module, the engineer:

\begin{enumerate}
  \item Declares range and modular types for register fields. Real-time
        diagnostics immediately flag any syntax errors.
  \item Defines the 32-bit GPIO control register struct with
        \texttt{[[offset/bit]]} annotations. The VSCode extension highlights
        the hardware attribute keywords distinctively.
  \item Declares the hardware register as a pointerless volatile variable
        at the base address.
  \item Writes an \texttt{init\_gpio} function with pre/post conditions.
        Go-to-definition resolves the register type.
  \item Runs \texttt{stcc gpio.stc} to generate SPARK Ada. The generated
        Ada includes record representation clauses, the \texttt{Address}
        aspect, and contract aspects directly usable by GNATprove.
  \item Runs GNATprove on the generated Ada to discharge the contracts.
\end{enumerate}

\begin{lstlisting}[caption={Complete GPIO driver module in stC}]
module gpio.driver {
  type modular 256    uint8_t;
  type range 0 .. 7  uint3_t;

  type struct {
    bool    enable     [[offset(0) bit(0)]];
    bool    irq_enable [[offset(0) bit(1)]];
    uint3_t mode       [[offset(0) bits(2..4)]];
    uint8_t pin_value  [[offset(2) bits(0..7)]];
    uint8_t irq_status [[offset(3) bits(0..7)]];
  } gpio_control_t;

  volatile gpio_control_t GPIO_A [[at(0x40000000)]];

  void gpio_init(in uint3_t initial_mode)
    [[pre(initial_mode <= 7)]]
    [[global(writes => (GPIO_A))]]
  {
    GPIO_A.enable     = true;
    GPIO_A.irq_enable = false;
    GPIO_A.mode       = initial_mode;
  }
}
\end{lstlisting}

The complete demo, including the generated SPARK Ada and a GNATprove run,
is available at \url{https://github.com/jgrivera67/stcc}.

%------------------------------------------------------------------
\section{Related Work}
\label{sec:related}
%------------------------------------------------------------------

\textbf{Unsafe C subsets and guidelines.} MISRA C~\cite{misra-c} and
CERT C define coding rules that restrict unsafe C constructs, but they remain
annotations on top of C and do not provide a formal type system or
machine-checkable contracts.

\textbf{Safe C alternatives.} Rust~\cite{rust} provides memory safety via
ownership types and has gained traction in embedded systems. Unlike stC, Rust
does not target SPARK verification and its ownership model, while powerful,
imposes a significant learning curve. Cforall~\cite{cforall} extends C with
parametric polymorphism but does not address formal verification.

\textbf{High-level synthesis languages.} SystemC and Bluespec target
hardware description and do not produce formally verifiable software.

\textbf{C to Ada transpilers.} No production transpiler from C to Ada exists;
the semantic gap is too large. stC avoids this by designing a new language
whose semantics are Ada from the ground up, with C surface syntax.

\textbf{SPARK subsets of Ada.} Spark2014~\cite{SparkAda} and GNAT Studio
provide IDE support for SPARK. stC complements these tools by offering a
less-verbose entry point for engineers unfamiliar with Ada syntax.

\textbf{Language server protocol tools.} Multiple academic languages
(e.g., F*~\cite{fstar}, Dafny~\cite{dafny}) ship VSCode extensions with
LSP support. stC follows this approach to lower the barrier to adoption.

%------------------------------------------------------------------
\section{Conclusion}
\label{sec:conclusion}
%------------------------------------------------------------------

stC demonstrates that the barrier between C-fluent embedded engineers
and formal verification need not be the syntax of the language they write in.
By providing a language that looks like C but compiles to SPARK Ada, stC
allows engineers to incrementally introduce formal contracts into existing
development workflows without abandoning familiar notation.

The tool is ready for demonstration on realistic embedded case studies.
Ongoing work focuses on: completing the stC-to-SPARK Ada code generator,
extending the type checker and contract verifier, adding array bounds
contracts, and packaging the toolchain as a single Alire crate.
We welcome contributions and industrial case studies from the FMICS community.

%------------------------------------------------------------------
\begin{thebibliography}{10}

\bibitem{SparkAda}
McCormick, J., Chapin, P.: Building High Integrity Applications with SPARK.
Cambridge University Press (2015)

\bibitem{misra-c}
MISRA Consortium: MISRA C:2012 -- Guidelines for the Use of the C Language
in Critical Systems. MIRA Ltd. (2012)

\bibitem{ada-industry}
Taft, S.T., Duff, R.A., Brukardt, R.L., Ploedereder, E., Leroy, P.:
Ada 2005 Reference Manual. LNCS 4348, Springer (2007)

\bibitem{alire}
Alire Project: Alire -- Ada/SPARK Package Manager.
\url{https://alire.ada.dev} (2024)

\bibitem{lsp}
Microsoft: Language Server Protocol Specification 3.17.
\url{https://microsoft.github.io/language-server-protocol/} (2023)

\bibitem{rust}
Klabnik, S., Nichols, C.: The Rust Programming Language.
No Starch Press (2019)

\bibitem{cforall}
Moss, P., et al.: C$\forall$: Adding Modern Programming Language Features to C.
SCP 161, 1--30 (2018)

\bibitem{fstar}
Swamy, N., et al.: Dependent Types and Multi-Monadic Effects in F*.
In: POPL 2016, ACM (2016)

\bibitem{dafny}
Leino, K.R.M.: Dafny: An Automatic Program Verifier for Functional Correctness.
In: LPAR-16. LNCS 6355, Springer (2010)

\end{thebibliography}

\end{document}
